\subsection{分次代数}

定义 1. 设 \(\Delta\) 为一个交换幺半群，A 为一个 \(\Delta\) 型分次环（II, § 11, 第 2 号）， \((\mathbf{A}_{\lambda})_{\lambda \in \Delta}\) 是它的分次，E 为一个 A-代数。加法群 E 上的一个分次 \((\mathbf{E}_{\lambda})_{\lambda \in \Delta}\) （型为 \(\Delta\) ）称为与 E 上的 A-代数结构相容，如果它既与 E 上的 A-模结构相容，又与 E 上的环结构相容；换言之，若对所有 \(\lambda, \mu\) （均属于 \(\Delta\) ），

\[
\mathrm{A} _ {\lambda} \mathrm{E} _ {\mu} \subset \mathrm{E} _ {\lambda + \mu}\tag{1}
\]

\[
\mathbf {E} _ {\lambda} \mathbf {E} _ {\mu} \subset \mathbf {E} _ {\lambda + \mu}\tag{2}
\]

A-代数 E 连同这个分次称为分次环 A 上的一个 \(\Delta\) 型分次代数。

当 A 上的分次为平凡分次（即（II, § 11, 第 1 号） \(\mathbf{A}_0 = \mathbf{A}, \mathbf{A}_{\lambda} = \{0\}\) ，对 \(\lambda \neq 0\) ）时，条件 (1) 意味着诸 \(\mathbf{E}_{\lambda}\) 是 E 的 A-子模。由此引出在非分次交换环 A 上定义 \(\Delta\) 型分次代数这一概念：给 A 赋予 \(\Delta\) 型的平凡分次，然后应用上述定义。

当我们考虑带单位元 \(e\) 的分次 A-代数 E 时，总约定 \(e\) 是 0 次的（参见练习 1）。

由此可知，若可逆元素 \(x \in E\) 是齐次的且次数为 p，则其逆 \(x^{-1}\) 也是齐次的且次数为 -p：只需把 \(x^{-1}\) 分解为关系 \(x^{-1}x = xx^{-1} = e\) 中的齐次元素之和即可 .

设 E 与 \(E'\) 是 \(\Delta\) 型分次环 A 上的两个 \(\Delta\) 型分次代数。A-代数同态 \(u: E \to E'\) 称为分次代数同态，如果 \(u(E_{\lambda}) \subset E_{\lambda}'\) 对所有 \(\lambda \in \Delta\) 成立（其中 \((E_{\lambda})\) 与 \((E_{\lambda}')\) 分别表示 E 与 \(E')\) 的分次）；当 E 与 \(E'\) 是结合的且含单位元而 u 保单位时，这个条件意味着 u 是分次环同态（II, §11, 第 2 号）。

设 E 是 N 型的 A-分次代数。通过令 \(E_{n} = \{0\}\) （对 n \textless{} 0）把 E 等同于 Z 型的 A-分次代数。

注. 定义 1 也可以这样解释：E 是一个分次 A-模，且 A-线性映射

\[
m \colon \mathbf {E} \otimes_ {\mathbf {A}} \mathbf {E} \rightarrow \mathbf {E}
\]

在 E 上定义乘法（§ 1, 第 3 号），当 E ⊗\_A E 被赋予其 Δ 型分次（II, § 11, 第 5 号）时，它是 0 次齐次的。

因此，在分次环 A 上定义一个以 E 为底层分次 A-模的 \(\Delta\) 型分次 A-代数结构，就相当于对每个有序对 \((\lambda, \mu)\) （由 \(\Delta\) 的元素组成）定义一个 Z-双线性映射

\[
m _ {\lambda \mu}: \mathrm{E} _ {\lambda} \times \mathrm{E} _ {\mu} \rightarrow \mathrm{E} _ {\lambda + \mu}
\]

使得对每个指标三元组 \((\lambda, \mu, \nu)\) 以及 \(\alpha \in \mathbf{A}_{\lambda}\) , \(x \in \mathbf{E}_{\mu}\) , \(y \in \mathbf{E}_{\nu}\) , \(\alpha \cdot m_{\mu\nu}(x, y) = m_{\lambda + \mu, \nu}(\alpha x, y) = m_{\mu, \lambda + \nu}(x, \alpha y)\) .

例。(1) 设 B 是 \(\Delta\) 型分次环；若赋予 B 其典范的 \(\mathbf{Z}\) -代数结构（ \(\S 1\) ，第 1 号，例 2），则 B 是分次 A-代数（ \(\mathbf{Z}\) 被赋予平凡分次）。

\begin{enumerate}
\def\labelenumi{(\arabic{enumi})}
\setcounter{enumi}{1}
\item
  设 \(\mathbf{A}\) 是 \(\Delta\) 型分次交换环， \(\mathbf{M}\) 是 \(\Delta\) 型分次 A-模。假设幺半群 \(\Delta\) 的所有元素都可消去，这使我们可以（II, § 11, 第 6 号）在 \(\operatorname{Homgr}_{\mathbf{A}}(\mathbf{M}, \mathbf{M}) = \operatorname{Endgr}_{\mathbf{A}}(\mathbf{M})\) 上定义一个 \(\Delta\) 型分次 A-模结构；由于这个分次与 \(\operatorname{Endgr}_{\mathbf{A}}(\mathbf{M})\) 上的环结构相容（II, § 11, 第 6 号），它就在 A-代数 \(\operatorname{Endgr}_{\mathbf{A}}(\mathbf{M})\) 上定义了一个含单位元的分次 A-代数结构 .
\item
  广群的代数。设 S 是一个广群， \(\phi: S \to \Delta\) 是一个同态。对所有 \(\lambda \in \Delta\) ，记 \(S_{\lambda} = \phi^{-1}(\lambda)\) ；于是 \(S_{\lambda}S_{\mu} \subset S_{\lambda + \mu}\) 。设 A 是 \(\Delta\) 型分次交换环， \((A_{\lambda})_{\lambda \in \Delta}\) 是它的分次；我们将在广群 S 的代数 \(E = A^{(S)}\) 上（§2, 第 6 号）定义一个分次 A-代数结构。为此，设 \(\mathbf{E}_{\lambda}\) 表示 \(\mathbf{E}\) 中由形如 \(\alpha .s\) （其中 \(\alpha \in \mathbf{A}_{\mu}, s \in \mathrm{S}_{\nu}\) 且 \(\mu + \nu = \lambda\) ）的元素生成的加法子群。由于诸 \(\mathrm{S}_{\lambda}\) 两两不相交， \(\mathbf{E}\) 是诸 \(\mathbf{A}_{\mu}\mathrm{S}_{\nu}\) 的直和，因而也是诸 \(\mathbf{E}_{\lambda}\) 的直和，并且显然诸 \(\mathbf{E}_{\lambda}\) 满足条件 (1) 与 (2)，从而在 \(\mathbf{E}\) 上定义了所要的分次 A-代数结构。若 \(\mathbf{S}\) 有单位元 \(e\) ，还可设 \(\phi(e) = 0\) 。特殊情形是环 \(\mathbf{A}\) 的分次为平凡的情形；此时 \(\mathbf{E}_{\lambda}\) 是 \(\mathbf{E}\) 的由 \(\mathrm{S}_{\lambda}\) 生成的 A-子模。更特殊地，若取 \(\mathbf{S} = \mathbf{N}^{(I)}, \Delta = \mathbf{N}\) ， \(\phi\) 为满足 \(\phi((n_{i})) = \sum_{i \in I} n_{i}\) 的映射，且环 \(\mathbf{A}\) 取平凡分次，便在多项式代数 \(\mathbf{A}[\mathbf{X}_{i}]_{i \in I}\) 上得到一个分次，对这个分次，齐次多项式 \(\neq 0\) 的次数就是 §2 第 9 号中定义的总次数（参见 §6 第 6 号）。
\end{enumerate}

现在取 S 为集合 B 的自由幺半群 Mo(B)（I, §7, 第 2 号）， \(\phi\) 为同态 \(\mathrm{Mo(B)} \to \mathbf{N}\) ，它把每个字映为其字长。这样就在集合 B 的自由结合代数上得到一个分次 A-代数结构（§ 2, 第 7 号；参见 § 5, 第 5 号）。

